\section{模型评价与结论}
\subsection{模型优点}
本文以同一能量平衡和SOC递推贯穿四问，新增机制只改变信息集和费用项，便于归因比较。预测按日期严格前推，避免随机划分造成时间泄漏；从1月1日开始连续推进SOC，避免2月1日人为重置状态；优化后独立回代6,570项约束并保存迭代记录。分位裕度将紧急购电代价转化为可解释的保守净负荷。
\subsection{模型局限}
分位鲁棒LP是完整场景风险模型的近似，未显式描述时序相关场景或输出CVaR。储能互斥依赖正电价、效率损耗和极小吞吐正则项，而非二进制变量；在允许负电价或售电时需改为混合整数模型。组合比较属于给定样本内回测，若用于长期运营，还应进行跨年度外推验证。
\subsection{结论与推广}
确定性条件下，储能低充高放使典型日费用下降26.90\%；严格历史预测与80\%残差分位裕度下，问题二334日费用为15,386,279.80元。固定价和波动价两种情形中，8种组合均由“不调整”获得最低累计费用。因此实际运营应把“是否调整”作为带成本的离散决策，而不是收到新预报后机械重算。该框架可推广至园区综合能源、充电站和含可再生能源的企业微网。
